\let\negmedspace\undefined
\let\negthickspace\undefined
\documentclass[journal]{IEEEtran}
\usepackage[a5paper, margin=10mm, onecolumn]{geometry}
\usepackage{lmodern} 
\usepackage{tfrupee} 
\setlength{\headheight}{1cm}
\setlength{\headsep}{0mm}   

\usepackage{gvv-book}
\usepackage{gvv}
\usepackage{cite}
\usepackage{amsmath,amssymb,amsfonts,amsthm}
\usepackage{algorithmic}
\usepackage{graphicx}
\usepackage{textcomp}
\usepackage{xcolor}
\usepackage{txfonts}
\usepackage{listings}
\usepackage{enumitem}
\usepackage{mathtools}
\usepackage{gensymb}
\usepackage{comment}
\usepackage[breaklinks=true]{hyperref}
\usepackage{tkz-euclide} 
\usepackage{listings}                             
\def\inputGnumericTable{}                                 
\usepackage[latin1]{inputenc}                                
\usepackage{color}                                            
\usepackage{array}                                            
\usepackage{longtable}                                       
\usepackage{calc}                                             
\usepackage{multirow}                                         
\usepackage{hhline}                                           
\usepackage{ifthen}                                           
\usepackage{lscape}
\usepackage{xparse}

\bibliographystyle{IEEEtran}

\title{5.2.64}
\author{EE25BTECH11062 - Vivek K Kumar}

\begin{document}
\maketitle

\renewcommand{\thefigure}{\theenumi}
\renewcommand{\thetable}{\theenumi}

\numberwithin{equation}{enumi}
\numberwithin{figure}{enumi} 

\textbf{Question}:\\
Solve for the following system of linear equations:
\begin{align*}
\vec{X} + \vec{Y} = \myvec{7 & 0 \\ 2 & 5}\\
\vec{X} - \vec{Y} = \myvec{3 & 0 \\ 0 & 4}\\
\end{align*}
\textbf{Solution: }
The given system can be represented as:
\begin{align}
    \myvec{\vec{I} & \vec{I} \\ \vec{I} & -\vec{I}}\myvec{\vec{X} \\ \vec{Y}} = \myvec{7 & 0 \\ 2 & 5 \\ 3& 0 \\ 0 & 4} 
\end{align}
Solving the equation,
\begin{align}
    \myvec{\vec{X} \\ \vec{Y}} = \myvec{1 & 0 & 1 & 0 \\ 0 & 1 & 0 & 1 \\ 1 & 0 & -1 & 0 \\ 0 & 1 & 0 & -1} ^{-1} \myvec{7 & 0 \\ 2 & 5 \\ 3& 0 \\ 0 & 4} \\
    \myvec{\vec{X} \\ \vec{Y}} = \frac{1}{2}\myvec{1 & 0 & 1 & 0 \\ 0 & 1 & 0 & 1 \\ 1 & 0 & -1 & 0 \\ 0 & 1 & 0 & -1} \myvec{7 & 0 \\ 2 & 5 \\ 3& 0 \\ 0 & 4} \\
    \myvec{\vec{X} \\ \vec{Y}} = \myvec{5 & 0 \\ 1 & 9/2 \\ 2 & 0 \\ 1 & 1/2}
\end{align}

Hence, the matrices are given by
\begin{align}
    \vec{X} = \myvec{5 & 0 \\ 1 & 9/2}\\
    \vec{Y} = \myvec{2 & 0 \\ 1 & 1/2}
\end{align}
\begin{table}[H]    
  \centering
  \begin{tabular}[12pt]{ |c| c|}
    \hline
    \textbf{Matrix} & \textbf{Value}\\ 
    \hline
    $\vec{X}$ & $\myvec{5 & 0 \\ 1 & 9/2}$\\
    \hline
    $\vec{Y}$ & $\myvec{2 & 0 \\ 1 & 1/2}$\\
    \hline
\end{tabular}
  \caption{Results}
  \label{tab:5.2.64}
\end{table}
\end{document}  